% Généré par CV Builder à partir de content.json et style.json.
\documentclass[10pt,a4paper]{article}

% -------------------- PACKAGES --------------------
\usepackage[margin=1cm]{geometry}
\usepackage{fontspec}
% Le texte extrait du PDF reprend exactement les caractères tapés (« - »,
% « · »…) : indispensable pour les ATS.
\XeTeXgenerateactualtext=1
\usepackage[french]{babel}
\frenchsetup{StandardLists=true}
\usepackage{enumitem}
\usepackage{hyperref}
\usepackage{xcolor}
\usepackage{graphicx}
\usepackage{tikz}

% -------------------- POLICE --------------------
% Ligatures « fi », « fl » coupées : un ATS lit alors des lettres normales.
\setmainfont{Lato}[
  Path=<fonts>/Lato/,
  Extension=.ttf,
  UprightFont=*-Regular,
  BoldFont=*-Bold,
  ItalicFont=*-Italic,
  BoldItalicFont=*-BoldItalic,
  Numbers=Proportional,
  Ligatures={TeX, CommonOff}]

% -------------------- COLORS --------------------
\definecolor{cvprimary}{HTML}{001F82}
\definecolor{cvsecondary}{HTML}{4F6B8A}
\definecolor{cvtext}{HTML}{000000}
\definecolor{cvaccent}{HTML}{2E7D32}
\definecolor{cvband}{HTML}{ECF3FB}
\definecolor{cvmuted}{HTML}{5B6472}
\definecolor{cvtitleband}{HTML}{E3EAF6}
\definecolor{cvlink}{HTML}{1F4FA3}

% -------------------- STYLE --------------------
\setlength{\parindent}{0pt}
\setlength{\parskip}{3pt}

% Titre de rubrique avec son filet. L'espace avant est explicite : dans une
% minipage, \section le perd.
\newcommand{\rubrique}[2][12pt plus 2fil]{%
  \par\vspace{#1}%
  {\fontsize{13}{15.5}\selectfont \textbf{\color{cvprimary}#2}\par}%
  \vspace{2pt}{\color{cvprimary}\hrule height 0.4pt}\vspace{7pt}}

\setlist[itemize]{label=\textbullet, leftmargin=12pt, labelsep=5pt, itemsep=1.5pt, parsep=0pt, topsep=2pt}
\pagestyle{empty}

% Espace entre deux entrées d'une rubrique, et entre deux blocs de compétences.
\newcommand{\entre}{\par\vspace{7pt plus 1fil}}
\newcommand{\petitentre}{\par\vspace{5pt plus 0.5fil}}
% Hauteur laissée aux deux colonnes : tout le reste de la page.
\newlength{\hauteurcol}


% -------------------- DOCUMENT --------------------
\begin{document}

% -------------------- HEADER --------------------
% La photo, carrée, a exactement la hauteur du bloc de texte à côté.
%@ header
\newsavebox{\blocentete}
\sbox{\blocentete}{\begin{minipage}[t]{0.78\textwidth}
    \vspace{0pt}\raggedright
    {\fontsize{24.88}{30}\selectfont \textbf{\color{cvprimary}CAMILLE LAURENT}}\par
    \vspace{12pt}
    {\fontsize{12}{14}\selectfont \textbf{\color{cvprimary}Élève en école d'ingénieurs | Data Science \& IA}}\par
    \vspace{7pt}
    En dernière année d'école d'ingénieurs, spécialité data science, à la recherche d'un stage de fin d'études en IA appliquée à partir de février 2027.\par
\end{minipage}}
\noindent\usebox{\blocentete}\hfill
\begin{minipage}[t]{\dimexpr\ht\blocentete+\dp\blocentete\relax}
    \vspace{0pt}
    \includegraphics[width=\linewidth]{imgcv.jpg}
\end{minipage}

\vspace{6pt}
{\color{cvprimary}\rule{\textwidth}{0.4pt}}
\vspace{2pt}

%@ contact 0
\textbf{Lyon, France} \quad | \quad \textbf{Téléphone :} 06 00 00 00 00 \quad | \quad \textbf{Email :} \href{mailto:camille.laurent@example.com}{camille.laurent@example.com} \quad | \quad \textbf{Permis B}

%@ contact 1
\textbf{LinkedIn :} \href{https://www.linkedin.com/in/camille-exemple/}{linkedin.com/in/camille-exemple} \quad | \quad \textbf{GitHub :} \href{https://github.com/camille-exemple}{github.com/camille-exemple} \quad | \quad \textbf{Mobilité :} toute la France

\vspace{10pt}

% -------------------- DEUX COLONNES --------------------
% Les deux colonnes prennent toute la hauteur restante ; l'espace libre se
% répartit entre les rubriques et les entrées.
%@ none
\par\setlength{\hauteurcol}{\dimexpr\pagegoal-\pagetotal-\baselineskip-4pt\relax}
\noindent
\begin{minipage}[t][\hauteurcol][s]{0.41\textwidth}
\raggedright
%@ none
% Fond de la colonne : du haut de la colonne au bas de la page, depuis le bord
% gauche de la page jusqu'au bout des filets. (Deux compilations.)
\begin{tikzpicture}[remember picture, overlay]
  \fill[cvband] ([yshift=14pt]current page.west |- 0,0) rectangle (\linewidth, 0 |- current page.south);
\end{tikzpicture}%

%@ section formation
\rubrique[0pt]{FORMATION}

%@ item formation 0
\textbf{Diplôme d'ingénieur, Data Science}\\
École d'ingénieurs du Numérique, Lyon\\
\textbf{\textit{2024 – 2027}}\\
3e année, option \textbf{apprentissage automatique} et statistique.

\entre
%@ item formation 1
\textbf{Classes préparatoires (MPSI, MP)}\\
Lycée du Parc (exemple), Lyon\\
\textbf{\textit{2022 – 2024}}\\
Mathématiques (algèbre linéaire, probabilités), physique et informatique.


%@ section competences
\rubrique{COMPÉTENCES}

%@ item competences 0
\textbf{Programmation}\\
Python, C/C++, Java, TypeScript, SQL.

\petitentre
%@ item competences 1
\textbf{Data \& IA}\\
Apprentissage automatique (scikit-learn, LightGBM), deep learning (PyTorch), vision par ordinateur, NLP, LLM, RAG, recherche sémantique.

\petitentre
%@ item competences 2
\textbf{Développement}\\
API REST (FastAPI), React, tests automatisés, architecture logicielle.

\petitentre
%@ item competences 3
\textbf{Outils}\\
Git, Docker, Linux, PostgreSQL, CI/CD, LaTeX.


%@ section soft-skills
\rubrique{SOFT SKILLS}

{\renewcommand{\arraystretch}{1.15}%
\begin{tabular}{@{}p{0.46\linewidth}@{\hspace{0.04\linewidth}}p{0.5\linewidth}@{}}
%@ item soft-skills 0
\textbullet\ Autonomie & \textbullet\ Travail en équipe \\
%@ item soft-skills 2
\textbullet\ Communication & \textbullet\ Adaptabilité \\
%@ item soft-skills 4
\multicolumn{2}{@{}l@{}}{\textbullet\ Gestion du temps et des priorités} \\
%@ item soft-skills 5
\multicolumn{2}{@{}l@{}}{\textbullet\ Esprit d'analyse et de synthèse} \\
\end{tabular}\par}


%@ section langues
\rubrique{LANGUES}

%@ item langues 0
\textbf{Français :} langue maternelle\\
%@ item langues 1
\textbf{Anglais :} courant (C1, TOEIC 945)\\
%@ item langues 2
\textbf{Espagnol :} intermédiaire (B1)\\
%@ item langues 3
\textbf{Italien :} notions


%@ section centres-interet
\rubrique{CENTRES D'INTÉRÊT}

\begin{itemize}[leftmargin=*]
%@ item centres-interet 0
    \item \textbf{Escalade} en club, deux séances par semaine
%@ item centres-interet 1
    \item \textbf{Photographie} de rue et retouche
%@ item centres-interet 2
    \item \textbf{Bénévolat} : soutien scolaire en mathématiques
%@ item centres-interet 3
    \item \textbf{Course à pied} : semi-marathon de Lyon
\end{itemize}


%@ none
\end{minipage}%
\hfill
\begin{minipage}[t][\hauteurcol][s]{0.55\textwidth}
\raggedright
%@ none
\setlist*[itemize]{itemsep=4pt plus 1pt, topsep=4pt plus 1pt}

%@ section experience
\rubrique[0pt]{EXPÉRIENCE PROFESSIONNELLE}

%@ item experience 0
{\fontsize{11}{13.5}\selectfont \textbf{Stage Data Scientist}\enspace{\fontsize{10}{12}\selectfont \textit{· Équipe de trois stagiaires}}\par}
{\fontsize{11}{13.5}\selectfont Logistica (exemple), Villeurbanne\hfill{\fontsize{10}{12}\selectfont \textbf{\textit{Juin 2026 – Août 2026}}}\par}\vspace{3pt}
\begin{itemize}
%@ bullet experience 0 0
    \item Prévision de la demande à 14 jours pour 1 200 références : modèles LightGBM sur l'historique des ventes, validation glissante dans le temps.
%@ bullet experience 0 1
    \item Pipeline de préparation des données en Python et SQL (PostgreSQL), relancé chaque nuit sous Docker.
%@ bullet experience 0 2
    \item API FastAPI qui sert les prévisions en production au logiciel de réapprovisionnement, avec tests automatisés.
%@ bullet experience 0 3
    \item Tableau de bord des écarts entre prévision et ventes, présenté chaque semaine à l'équipe achats.
\end{itemize}

\entre
%@ item experience 1
{\fontsize{11}{13.5}\selectfont \textbf{Conseil en librairie (job étudiant)}\enspace{\fontsize{10}{12}\selectfont \textit{}}\par}
{\fontsize{11}{13.5}\selectfont Librairie (exemple), Lyon\hfill{\fontsize{10}{12}\selectfont \textbf{\textit{Septembre 2025 – Aujourd'hui}}}\par}\vspace{3pt}
\begin{itemize}
%@ bullet experience 1 0
    \item Accueil et conseil des clients, commandes spéciales et suivi des livraisons.
%@ bullet experience 1 1
    \item Inventaire mensuel d'un rayon et propositions de réassort à partir des ventes.
\end{itemize}


%@ section projets
\rubrique{PROJETS}

%@ item projets 0
\par\textbf{TrafficLens – Comptage de véhicules}\enspace{\fontsize{9}{11}\selectfont \textit{· Projet d'école}}\hfill\textbf{\textit{2026}}\par
Détection et suivi de véhicules sur vidéo (YOLO, PyTorch), comptage par voie et par heure.
\entre
%@ item projets 1
\par\textbf{Assistant de questions sur des documents}\enspace{\fontsize{9}{11}\selectfont \textit{· Projet personnel}}\hfill\textbf{\textit{2026}}\par
Recherche sémantique et réponses sourcées sur des PDF (embeddings, pgvector, FastAPI, React).
\entre
%@ item projets 2
\par\textbf{Planificateur de tournées}\enspace{\fontsize{9}{11}\selectfont \textit{· En binôme}}\hfill\textbf{\textit{2025}}\par
Optimisation de tournées de livraison (heuristiques, OR-Tools) et carte interactive (React).
\entre
%@ item projets 3
\par\textbf{Simulation de fluides 2D}\hfill\textbf{\textit{2024}}\par
Différences finies en C++ et rendu temps réel avec OpenGL.


%@ section associatif
\rubrique{TRAVAIL ASSOCIATIF}

%@ item associatif 0
\textbf{Trésorerie du bureau des sports}\\
Association sportive de l'école \hfill \textbf{\textit{2024 – Aujourd'hui}}\par\vspace{3pt}
\begin{itemize}
%@ bullet associatif 0 0
    \item Budget annuel des activités sportives, suivi des dépenses et des subventions.
%@ bullet associatif 0 1
    \item Organisation d'un tournoi inter-écoles de 200 participants.
\end{itemize}


%@ none
\end{minipage}

\end{document}
